\chapter{不完全分解}\label{v0ch:8}
\begin{chaptersummary}
精确稀疏分解可能产生昂贵的填充，因此求解目标转向构造低成本的近似逆作用。本章介绍不完全分解如何选择和丢弃条目，并分析主元与三角求解带来的约束。它为第四部分的预条件设计建立基础：近似的价值由外层迭代改善与构造、应用成本共同衡量。
\end{chaptersummary}
\learninggoals{区分 ILU0、ILUk、ILUT 与 IC；解释丢弃和填充；把因子作为解三角系统的操作。}
\paragraph{核心问题}哪些消元信息值得保留，才能用有限内存改善迭代？

\section{截断消元}
精确 LU 的更新可能在原零位置产生条目。ILU(0) 只保留预定的原始非零模式；ILU($k$) 根据符号消元的填充层级保留条目，原始条目层级为零，经主元 $j$ 产生的新条目层级取相关路径层级之和加一，并保留最小层级。ILUT 主要用数值阈值和每行条目限制控制保留量。不同库的阈值缩放、排序和主元选项需查其定义。
获得近似 $A\approx LU=M$ 后，预条件作用是
\[
 \bm z=M^{-1}\bm r:\qquad L\bm y=\bm r,\quad U\bm z=\bm y.
\]
通常不形成显式逆矩阵。
\section{最小可视算例}
三对角 Poisson 无主元 LU 不产生带外填充，因此该例中 ILU(0) 可等于精确 LU。四节点环状图或二维五点矩阵则会出现额外填充，ILU(0) 将其中一部分丢弃，得到真正的近似分解。先比较这两个例子，才能观察丢弃机制。
IC 在对称模式下近似 $A\approx LL^T$，用于 SPD 矩阵 $A$。丢弃条目会改变消元中的 Schur 补，因此还需检查 IC 主元的正性；遇到非正主元时，可调整排序、加入位移或采用改进策略。
\conceptbridge{近似分解既然不精确，为什么还值得计算？}{外层迭代只需利用一个便宜的修正方向变换，不要求预条件器独立完成求解。保留主要局部耦合可降低外层迭代成本；代价是每次应用要执行有依赖关系的三角求解。}{精确消元成本 $\longrightarrow$ 受控截断 $\longrightarrow$ 外层迭代中的近似逆。}
\costnote{更小丢弃阈值常增大填充、构造和三角求解成本，却不保证总时间单调改善。通用 ILU 通常不对称，不能直接作为标准 PCG 的 SPD 预条件器。零主元、极小主元和并行依赖是独立风险。}
\coderef{sparse_linear_algebra/solvers.py}
\aiexercise{分别对三对角与二维五点 Poisson 执行 ILU(0)，列出被丢弃的条目并比较预条件后的真实残差。}{扫描 ILUT 丢弃阈值，记录构造时间、因子非零数、GMRES 次数及总时间，寻找成本折中。}
\studyroute{略去修改不完全分解及其谱等价证明。参考 \cite{saad,templates}。}{ILU(0); level of fill; ILUT; incomplete Cholesky; drop tolerance}
